\subsection{Trayectoria, corte y transferencia}

Los pares excluidos son $O(n^3)$ y cada contribución ideal es $O(n^{-2})$, de modo que su efecto
normalizado es $O(n^{-1})$. Las cuadraturas y el límite del perfil dan, para $y=T/N$,
\begin{multline*}
 \frac{\Phi_0}{n^2}\longrightarrow b_w(y)=y^2\int w(r,c)\,\nu_1(r,j)\,\nu_1(i,c)\\
 \times\bigl[1-y(1-u_{D}(i-j))\bigr]\bigl[1-y(1-u_{A}(i+j))\bigr]\,dr\,dc\,di\,dj ,
\end{multline*}
uniformemente en $y$ en compactos, con $\nu_1$ el perfil lleno (de masa uno).

\begin{proposicion}[Cuenta inversa hasta el final]\label{prop:inversa}
Sea $0\le h\le K$ y $w$ continua y acotada. Para toda tolerancia fija $\gamma>0$,
\[
 \QQ_h\Bigl(S\cap\Bigl\{\Bigl|\frac{I_w}{n^2}-B_w\Bigr|>\gamma\Bigr\}\Bigr)\le e^{-c_{\gamma,w}n},
 \qquad B_w=\int w(r,c)\,\nu_1(r,j)\,\nu_1(i,c)\,\psi(i,j)\,dr\,dc\,di\,dj,
\]
donde $I_w$ se evalúa en la colocación final. En consecuencia, bajo $\mu_h$,
$I_w/n^2\to B_w$ en probabilidad y en media.
\end{proposicion}

\begin{proof}
Se corta cuando faltan $m_\delta=\lceil\delta n\rceil$ piezas reales, con $\delta\in(0,1/4)$ libre (no el $\delta_0$ de la Sección~\ref{sec:absorbedores}). El control de los colores,
$q_t=qz_t+O(\eps n)$ con $q/n\to1$, asegura $\tau_\delta\le T_{\delta/4}$ e
$y_{\tau_\delta}=1-\delta+O(\eps+n^{-1})$; se fija $\delta$, después $\eps\ll\delta$ y después
$n$. La unión sobre horizontes deterministas permite el corte aleatorio. Una reina nueva crea a
lo sumo $2n$ pares inversos y destruye a lo sumo $2n$ por sus dos diagonales; las piezas
virtuales no cambian $I_w$. Así, los últimos $m_\delta$ pasos reales cambian $I_w$ en a lo sumo
$4nm_\delta\|w\|_\infty$. Para cada tolerancia final fija se escoge $\delta$ suficientemente
pequeño, luego las tolerancias de~\eqref{eq:regularidad} y finalmente $n$ grande; por continuidad de
$b_w$ en $y=1$ y $b_w(1)=B_w$, el Lema~\ref{lem:conc_inversa} da la cota bajo $\QQ_h$.

Dividiendo por $\QQ_h(S)\ge e^{-d_n}$ se conserva una cota $e^{-cn+o(n)}$ bajo $\nu'_h$; el
Teorema~\ref{thm:comparacion} y el Lema~\ref{lem:binaria} dan la convergencia en probabilidad bajo
$\mu_h$, y la cota determinista $|I_w|/n^2\le\|w\|_\infty$ da la convergencia en media.
\end{proof}

No se afirma una cola exponencial bajo $\mu_h$ ni una tasa para el error de su esperanza.
